\documentclass[notes=show]{beamer}%
\usepackage{color}
\usepackage{mathpazo}
\usepackage{hyperref}
\usepackage{multimedia}
\usepackage{amsmath}
\usepackage{amsfonts}
\usepackage{amssymb}
\usepackage{graphicx}
\usepackage{subcaption} 
\usepackage{pgfplots}
\usepackage{outlines}
\usepackage{remreset}%
\setcounter{MaxMatrixCols}{30}
\useoutertheme{miniframes}
\defbeamertemplate{mini frame}{scaled circle}[1]
{%
  \begin{pgfpicture}{0pt}{0pt}{#1 * 0.1cm}{#1 * 0.1cm}
    \pgfsetlinewidth{#1 * 0.4pt}
    \pgfpathcircle{\pgfpoint{#1 * 0.05cm}{#1 * 0.05cm}}{#1 * 0.05cm}
    \pgfusepath{fill,stroke}
  \end{pgfpicture}%
}
[action]
{
  \newlength{\myminiframesize}
  \setlength{\myminiframesize}{0.14cm}
  \newlength{\myminiframeoffset}
  \setlength{\myminiframeoffset}{0.03cm}
  \setbeamersize{mini frame size=#1\myminiframesize,mini frame offset=#1\myminiframeoffset}
}

\defbeamertemplate{mini frame in current section}{scaled circle}[1]
{%
  \begin{pgfpicture}{0pt}{0pt}{#1 * 0.1cm}{#1 * 0.1cm}
    \pgfsetlinewidth{#1 * 0.4pt}
    \pgfpathcircle{\pgfpoint{#1 * 0.05cm}{#1 * 0.05cm}}{#1 * 0.05cm}
    \pgfusepath{stroke}
  \end{pgfpicture}%
}

\defbeamertemplate{mini frame in current subsection}{scaled circle}[1]
{%
  \begin{pgfpicture}{0pt}{0pt}{#1 * 0.1cm}{#1 * 0.1cm}
    \pgfsetlinewidth{#1 * 0.4pt}
    \pgfpathcircle{\pgfpoint{#1 * 0.05cm}{#1 * 0.05cm}}{#1 * 0.05cm}
    \pgfusepath{stroke}
  \end{pgfpicture}%
}

\setbeamertemplate{mini frame}[scaled circle]{0.6}
\setbeamertemplate{mini frame in current section}[scaled circle]{0.6}
\setbeamertemplate{mini frame in current subsection}[scaled circle]{0.6}

%TCIDATA{OutputFilter=latex2.dll}
%TCIDATA{Version=5.50.0.2953}
%TCIDATA{CSTFile=beamer.cst}
%TCIDATA{Created=Wednesday, December 14, 2016 08:32:20}
%TCIDATA{LastRevised=Monday, October 02, 2017 09:01:02}
%TCIDATA{<META NAME="GraphicsSave" CONTENT="32">}
%TCIDATA{<META NAME="SaveForMode" CONTENT="1">}
%TCIDATA{BibliographyScheme=Manual}
%TCIDATA{<META NAME="DocumentShell" CONTENT="Other Documents\SW\Slides - Beamer">}
%TCIDATA{ComputeDefs=
%1$W_{e}^{h}(z)=I_{e}+z+\max\limits_{\hat{z}}\{-\hat{z}+\beta\alpha
%_{h}[\upsilon(q)-\rho d]+\beta\rho\hat{z}+\delta W_{0}^{h}\left(  0\right)
%+(1-\delta)W_{1}^{h}\left(  0\right)  $
%}
%BeginMSIPreambleData
\providecommand{\U}[1]{\protect\rule{.1in}{.1in}}
%EndMSIPreambleData
\makeatletter
\@removefromreset{subsection}{section}
\let\beamer@writeslidentry@miniframeson=\beamer@writeslidentry
\def\beamer@writeslidentry@miniframesoff{  \expandafter\beamer@ifempty\expandafter{\beamer@framestartpage}{}  {        \clearpage\beamer@notesactions  }
}
\newcommand*{\miniframeson}
{\let\beamer@writeslidentry=\beamer@writeslidentry@miniframeson}
\newcommand*{\miniframesoff}
{\let\beamer@writeslidentry=\beamer@writeslidentry@miniframesoff}
\makeatother
\setcounter{subsection}{1}
\addtobeamertemplate{navigation symbols}{}{    \usebeamerfont{footline}    \usebeamercolor[fg]{footline}    \hspace{1em}    \insertframenumber/\inserttotalframenumber
}
\setbeamercolor{footline}{fg=blue}
\setbeamerfont{footline}{series=\bfseries}
\newenvironment{stepenumerate}{\begin{enumerate}[<+->]}{\end{enumerate}}
\newenvironment{stepitemize}{\begin{itemize}[<+->]}{\end{itemize} }
\newenvironment{stepenumeratewithalert}{\begin{enumerate}[<+-| alert@+>]}{\end{enumerate}}
\newenvironment{stepitemizewithalert}{\begin{itemize}[<+-| alert@+>]}{\end{itemize} }
\usetheme{Frankfurt}
\miniframesoff
\AtBeginSection[]{
\begin{frame}
\vfill
\centering
\begin{beamercolorbox}[sep=8pt,center,shadow=true,rounded=true]{title}
\usebeamerfont{title}\insertsectionhead\par  \end{beamercolorbox}
\vfill
\end{frame}
}
\providecommand{\U}[1]{\protect\rule{.1in}{.1in}}
\AtBeginSection[]{
\begin{frame}
\vfill
\centering
\begin{beamercolorbox}[sep=8pt,center,shadow=true,rounded=true]{title}
\usebeamerfont{title}\secname\par  \end{beamercolorbox}
\vfill
\end{frame}
}
\newtheorem{proposition}[theorem]{Proposition}
\newtheorem{remarks}[theorem]{Remarks}
\newtheorem{remark}[theorem]{Remark}
\newtheorem{conjecture}[theorem]{Conjecture}

\begin{document}

\title{Kalman Filter}
%\author[shortname]{Gauthier Vermandel}
%\institute[shortinst]{Universit\'e Paris Dauphine -- PSL Research}


\date{\today}
\maketitle

\section{Introduction}

\miniframeson

\begin{frame}
\frametitle{Parametrization}


\noindent There are different methods to parametrize models:\\
\bigskip
\begin{center}
\begin{itemize}

\item Calibration
\item Matching Moments
\item Estimation
 \begin{itemize}
\item Maximum Likelihood
\item Bayesian Estimation
\end{itemize}
\end{itemize}

\end{center}

\end{frame}

\begin{frame}
\frametitle{Calibration}
In a nutshell:
\begin{center}
 \begin{itemize}
\item Giving values to structural parameters in such way to replicate real world observations,
\item Use different values to perform sensitivity analysis.
\end{itemize}
\end{center}
\bigskip
The problem we calibration:
\begin{itemize}
\item Difficult to chose the right values for the parameters,
\item I could lead (sometimes) to arbitrary values with very little economic sens. 
\end{itemize}
\end{frame}

\begin{frame}

\frametitle{Matching Moments}

Similarly to calibration:
\begin{center}
 \begin{itemize}
\item Using values for parameters by choosing a set of moments that we would like our model to match,
\item Performing a sensitivity analysis to judge the model fit by investigating the set of moments.
\end{itemize}
\end{center}
\bigskip
Contrary to calibration:
\begin{itemize}
\item estimated (different than the estimation we will use for Kalman filter)
\item Allow for hypothesis testing 
\end{itemize}
\end{frame}


\begin{frame}

\frametitle{Estimation}

Full information method as we have to specify an entire distribution:
\begin{center}
 \begin{itemize}
\item \textbf{Maximum Likelihood:} The model fit is then based on implied likelihood function (i.e. how consistent with the distribution assumption we have made).
\item \textbf{Bayesian estimation:}  Need to specify a prior (will be discussed further on Bayesian estimation techniques).
\end{itemize}
\end{center}

\end{frame}


\miniframesoff

\section{Maximum Likelihood}

\miniframeson

\begin{frame}
\frametitle{Kalman Filter}

The Kalman filter was initially developed for engineering control problems (e.g. NASA spatial program)

\begin{center}
 \begin{itemize}
\item Following Hamilton notation we will set up the Kalman filter based on general time series model,
\item Then the time series model will be linked to DSGE models.
\end{itemize}
\end{center}

\end{frame}


\begin{frame}
\frametitle{Kalman Filter: Times Series Model}
The model:
\bigskip
\begin{equation}
y_{t}= M_{t}' x_{t}+w_{t}, \text{     } E(w_{t},w_{t}')=R, \label{obs eq}%
\end{equation}
\begin{equation}
x_{t+1}= F_{t} x_{t}+v_{t+1}, \text{     } E(v_{t},v_{t}')=Q.\label{state eq}%
\end{equation}

\begin{center}
 \begin{itemize}
\item $y_{t}$ is an observed variable, while $x_{t}$ is not observed,
\end{itemize}
\end{center}

$\Rightarrow$ Kalman filter allows us to retrieve an estimate of $x_{t}$.

\end{frame}


\begin{frame}
\frametitle{Kalman Filter: Some initial hypotheses}

\begin{center}
 \begin{itemize}
\item we assume that $w$ and $v$ are uncorrelated (we can easily expend the model to account for cross correlations), 
\item we assume for now that $H$, $F$, $R$, and $Q$ are known and time unvarying (without a loss of generality). (Later on we will show how to estimate them),
\item we need to specify an initial condition for $x_{t}$ as it is a forward looking variable: $x_{1}$ with mean $\hat{x}_{1}$ and variance $P_{1|0}$,
\item observation ($w_{t}$) and state ($v_{t}$) disturbances are:
 \begin{itemize}
\item uncorrelated over time,
\item uncorrelated with each other (at all leads and lags)
\item orthogonal to $x_{1}$
\end{itemize}
\end{itemize}
\end{center}

\end{frame}


\begin{frame}
\frametitle{Deriving the Kalman Filter}


The main idea of Kalman filter is to give a prediction about a variable of interest (e.g. potential GDP) giving that we observed a series of a variable values in the past (e.g. GDP):
\begin{equation}
\hat{x}_{t+1}= \hat{E}(x_{t+1}|Y_{t}), \text{ where:  } Y_{t}=(y_{t}',y_{t-1}',....,y_{1}').
\end{equation}
Now the question is how well our prediction? 
\begin{itemize}
\item Thus, a need to specify a loss function to be able to evaluate the forecast:
\end{itemize}
\begin{equation}
minE(x_{t+1}-\hat{x}_{t+1|t})^{2} \label{loss function eq}%
\end{equation}
\begin{itemize}
\item $\hat{x}_{t+1|t}$ is then the projection of $x_{t+1}$ on its regressors.
\end{itemize}

\end{frame}


\begin{frame}
\frametitle{Deriving the Kalman Filter}

\begin{itemize}
\item Let's recall that any projection of $x$ on $Y$ is a function of $\hat{E}[x|Y]=f(x)$. 
\item Assuming then both a linear form for the projection where $f(x)=\alpha+\beta x$ and the existence of both first and second moments\footnote{$\bar{x}=E(x)$, $\text{Cov}(x,x)=E[(x-\bar{x})(x-\bar{x})']$, $\bar{Y}=E(Y)$, $\text{Cov}(Y,Y)=E[(Y-\bar{Y})(Y-\bar{Y})']$, and $\text{Cov}(x,Y)=E[(x-\bar{x})(Y-\bar{Y})']$.}
\end{itemize}

We then solve the minimization problem (in the general framework before going back to our time series model) as follow: 

\begin{align}
\beta&=\frac{Cov(x,Y)}{Cov(Y,Y)},\\
\alpha&=\bar{x}-\beta \bar{Y},\\
\hat{x}&=\hat{E}(x|Y)=\bar{x}+\frac{Cov(x,Y)}{Cov(Y,Y)}(Y-\bar{Y}).
\end{align}

\end{frame}


\begin{frame}
\frametitle{Deriving the Kalman Filter}

\textbf{We should not confuse projection with regression (e.g. OLS). Although the expression we are seeking to minimize are closely similar, OLS seeks to investigate a causality, while the projection is solely focused on forecasting.}
\end{frame}


\begin{frame}
\frametitle{Deriving the Kalman Filter}

Jumping back to our model! 

The main idea is to write down the model recursively:

\begin{itemize}
\item As we have a value for $x_{1|0}$, and the observation for $Y_{t}$ (i.e. $y_{1}$,....,$y_{T}$), we can then have an update for $\hat{x}_{1|1}$ and then have a forecast value for $\hat{x}_{2|1}$. Then we keep repeating this procedure.
\end{itemize}

Thus,

\begin{align}
\hat{x}_{t|t}&=\hat{E}(x_{t}|Y_{t}) \nonumber \\
&=\hat{x}_{t|t-1}+E[(x_{t}-\hat{x}_{t|t-1})(y_{t}-\hat{y}_{t|t-1})']* \nonumber\\
&E[(y_{t}-\hat{y}_{t|t-1})(y_{t}-\hat{y}_{t|t-1})']^{-1}(y_{t}-\hat{y}_{t|t-1})).
\end{align}

\end{frame}


\begin{frame}
\frametitle{Deriving the Kalman Filter}

Using the model equations, we can easily rewrite and simplify the co-variance, the variance, and error terms, respectively, as following:

\begin{align}
E[(x_{t}-\hat{x}_{t|t-1})(y_{t}-\hat{y}_{t|t-1})']&=E[(x_{t}-\hat{x}_{t|t-1})(M'(x_{t}-\hat{x}_{t|t-1})+w_{t})']\nonumber\\
&=E[(x_{t}-\hat{x}_{t|t-1})((x_{t}-\hat{x}_{t|t-1})'M]\nonumber\\
&=P_{t|t-1}M
\end{align}
\begin{align}
E[(y_{t}-\hat{y}_{t|t-1})(y_{t}-\hat{y}_{t|t-1})']&=E[M'(x_{t}-\hat{x}_{t|t-1})(x_{t}-\hat{x}_{t|t-1})M]\nonumber\\
&+E[w_{t}w{t}']\nonumber\\
&=M'P_{t|t-1}M+R
\end{align}

\begin{align}
y_{t}-\hat{y}_{t|t-1} &=y_{t} - M'\hat{x}_{t|t-1}
\end{align}

\end{frame}

\begin{frame}
\frametitle{writing down the final form of the Kalman Filter}
\underline{The state variable:} \\
\bigskip
\textbf{The update step:} 
\begin{align}
\hat{x}_{t|t}=\hat{x}_{t|t-1}+P_{t|t-1}M(M'P_{t|t-1}M+R)^{-1}(y_{t}-M'\hat{x}_{t|t-1})
\end{align}
\textbf{The forecast step:} 
\begin{align}
\hat{x}_{t+1|t}&=\hat{E}(x_{t+1}|Y_{t})\nonumber\\
&=F\hat{E}(x_{t}|Y_{t})\nonumber\\
&=F\hat{x}_{t|t-1}+FP_{t|t-1}M(M'P_{t|t-1}M+R)^{-1}(y_{t}-M'\hat{x}_{t|t-1})
\end{align}

\end{frame}

\begin{frame}
\frametitle{writing down the final form of the Kalman Filter}
\underline{Specifying the variance matrix recursion formula:} \\
\bigskip
\textbf{The update step:} 
\begin{align}
P_{t|t}&=E[(x_{t}-\hat{x}_{t|t})(x_{t}-\hat{x}_{t|t})']\nonumber\\
&=E[(x_{t}-\hat{x}_{t|t-1})(x_{t}-\hat{x}_{t|t-1})'] - E[(x_{t}-\hat{x}_{t|t-1})(y_{t}-\hat{y}_{t|t-1})']*\nonumber\\
&E[(y_{t}-\hat{y}_{t|t-1})(y_{t}-\hat{y}_{t|t-1})']^{-1}*\nonumber\\
&E[(y_{t}-\hat{y}_{t|t-1})(x_{t}-\hat{x}_{t|t-1})']\nonumber\\
&=P_{t|t-1}-P_{t|t-1}M(M'P_{t|t-1}M+R)^{-1}M'P_{t|t-1}
\end{align}
\textbf{The forecast step:} 
\begin{align}
P_{t+1|t}&=E[(x_{t+1}-\hat{x}_{t+1|t})(x_{t+1}-\hat{x}_{t+1|t})']\nonumber\\
&=E[(F x_{t}+v_{t+1}- F\hat{x}_{t|t})(F x_{t}+v_{t+1}-F\hat{x}_{t|t})']\nonumber\\
&=F E[(x_{t}-\hat{x}_{t|t})(x_{t}-\hat{x}_{t|t})']F'+E[v_{t+1}v_{t+1}']\nonumber\\
&=FP_{t|t}F'+Q
\end{align}

\end{frame}


\begin{frame}
\frametitle{Summary of the Kalman Filter}
\textbf{The update step:} 
\begin{align}
&\hat{x}_{t|t}=\hat{x}_{t|t-1}+P_{t|t-1}M(M'P_{t|t-1}M+R)^{-1}(y_{t}-M'\hat{x}_{t|t-1})\\
&P_{t|t}=P_{t|t-1}-P_{t|t-1}M(M'P_{t|t-1}M+R)^{-1}M'P_{t|t-1}
\end{align}
\textbf{The forecast step:} 
\begin{align}
&\hat{x}_{t+1|t}=F\hat{x}_{t|t-1}+FP_{t|t-1}M(M'P_{t|t-1}M+R)^{-1}(y_{t}-M'\hat{x}_{t|t-1})\\
&P_{t+1|t}=FP_{t|t}F'+Q
\end{align}

\end{frame}

\miniframesoff
\section{Maximum Likelihood (ML)}

\miniframeson

\begin{frame}
\frametitle{Estimation}

\begin{itemize}
\item We can also estimate all the parameter we assumed previously fixed (i.e. $\varkappa=\{M, F, Q, R\}$) using ML,
\item Assuming $x_{1|0}$ and both error terms are Gaussian, then the distribution of $y_{t}/Y_{t}$ is also Gaussian
\item We then given values of $\varkappa$ we can derive the mean and variance of y,
\item And as we know the distribution of y, this allows for computing the likelihood of the observations $y_{t}$
\item Finally, we just maximize the log-likelihood with respect to the parameters $\varkappa$
\end{itemize}


\end{frame}

\miniframesoff
\section{Liking the Time Series Model to DSGE}

\miniframeson

\begin{frame}

\frametitle{DSGE Example: Neo-classical Growth Model}
\textbf{Household:} 
The household maximize their lifetime utility. They consume and invest in capital.

\begin{align}
\max_{\left\{  c_{t},k_{t+1}\right\}  }E_{0}\sum_{t=0}^{\infty
}\beta^{t}\left( log(c_{t})\right),\nonumber\\
c_{t}+k_{t+1}=y_{t}+(1-\delta)k_{t}
\end{align}
(with initial value for $k_{0}$ known.)

\textbf{Production:} 
The production follow a classical cobb-douglas function and is subject to exogenous technology shocks
\begin{align}
y_{t}&=a_{t}k_{t}^\alpha\\
a_{t}&=1-\rho+\rho a_{t-1}+\epsilon_{t}
\end{align}
with $\epsilon_{t}$$\sim$ N(0,$\sigma^2$) and $a_{0}$ known.


\end{frame}


\begin{frame}

\frametitle{DSGE Example: Neo-classical Growth Model}
\textbf{Solving the modeling (i.e. finding the policy functions)} \\
\bigskip
The first order conditions read:

\begin{align}
c_{t}^{-1}&=E_{t}[\beta c_{t+1}^{-1}(\alpha z_{t+1}k_{t}^(\alpha-1)+1-\delta)]\\
c_{t}+k_{t}&=a_{t}k_{t-1}^\alpha+(1-\delta)k_{t-1}
\end{align}

Thus, the policy function can be written as following:
\begin{align}
k_{t}&=k_{ss}+\alpha_{kk}(k_{t-1}-k_{ss})+ \alpha_{ka}(a_{t}-a_{ss})\\
a_{t}&=(1-\rho)a_{ss}+\rho a_{t-1}+\epsilon_{t} 
\end{align}
where $a_{ss}=1$.

\begin{itemize}
\item $\alpha_{kk}$, $\alpha_{ka}$, $k_{ss}$ are functions which depend on the structural parameters $\varkappa=\{\alpha,\beta,\delta,\rho,\sigma,a_{0}\}$.
\end{itemize}

\end{frame}

\begin{frame}

\frametitle{DSGE Example: Neo-classical Growth Model}
\textbf{Estimating the model: Using the state-space form} \\
\bigskip
\begin{itemize}
\item Using ML we can estimate the model
\item To do so we need to specify the likelihood function, thus the Kalman filter M and F matrices equivalents.
\end{itemize}
\bigskip
\textbf{The state-space form (let's recall the general case)} \\

\begin{align}
y_{t}= M_{t}' x_{t}+w_{t}, \text{     } E(w_{t},w_{t}')=R, \nonumber\\
x_{t+1}= F_{t} x_{t}+v_{t+1}, \text{     } E(v_{t},v_{t}')=Q,\nonumber
\end{align}


\end{frame}

\begin{frame}

\frametitle{DSGE Example: Neo-classical Growth Model}
\textbf{Estimating the model: Using the state-space form} \\
\bigskip
Our Neo-classical Growth model:

\begin{equation*}
\begin{bmatrix}
k_{t} - k_{ss} \\
a_{t+1}-a_{ss} 
\end{bmatrix}	
 = 
\begin{bmatrix}
\alpha_{kk} & \alpha_{ka} \\
0 & \rho  
\end{bmatrix}	
\begin{bmatrix}
k_{t-1} - k_{ss}\\
a_{t}-a_{ss}   
\end{bmatrix}	
+
\begin{bmatrix}
0\\
\epsilon_{t+1}   
\end{bmatrix}
\end{equation*}


\begin{equation*}
k_{t} - k_{ss}
= 
\begin{bmatrix}
1 & 0  
\end{bmatrix}	
\begin{bmatrix}
k_{t-1} - k_{ss}\\
a_{t}-a_{ss}   
\end{bmatrix}	
+
\begin{bmatrix}
0
\end{bmatrix}
\end{equation*}

The work is done, now we just use the previous section results to compute the Kalman filter, while using ML through maximizing the log-likelihood function with respect to the parameters $\varkappa$.

\end{frame}

\begin{frame}

\frametitle{DSGE Example: Neo-classical Growth Model}
\textbf{Estimating the model: Using the state-space form} \\

The log-likelihood function for the Gaussian distribution assumed in the beginning reads as following:
\begin{align}
\mathcal{L}(Y_{t}|\varkappa) = log (f(y_{0}))+ \sum_{t=1}^{T} log (f(y_{t}|Y_{t-1};\varkappa))
\end{align}
which in this example would be:
\begin{align}
log\left(\frac{1}{\sigma \sqrt{2\pi}}\right) - \frac{\epsilon(\varkappa)^2}{2\sigma^2}
\end{align}
\end{frame}

\miniframesoff
\section{Exercise}

\miniframeson

\begin{frame}

\frametitle{Constructing Kalman Filter}

We will be using Kalman filter to estimate the aggregate matching function using ML.\\
\bigskip

Using the aggregate matching function, which allows for capturing the relationship between the number of new jobs (hires) $H_{t}$, and the number of unemployed $U_{t}$ as well as the number of vacancies present in the economy $V_{t}$, is a wide spread technique to measure unemployment. We characteristic the functional form following a classical Cobb Douglas constant returns to scale function.

\begin{align}
H_{t} = M U_{t}^{1-\rho}V_{t}^{\rho}
\end{align}
\begin{itemize}

\item where, $M$ the constant level parameter referring to the degree of mismatch or matching efficiency,
\item $\rho$ the elasticity of matches with respect to the number of unemployed,
\end{itemize}

\end{frame}


\begin{frame}

\frametitle{Constructing Kalman Filter}

The goal of this exercise is to attempt an estimation of the aggregate matching function, while allowing match efficiency to 
fluctuate over time, in order to investigate to role of the mismatch in the unemployment spike
during the Great Recession. 

\end{frame}


\begin{frame}

\frametitle{First Step:}
\textbf{Estimating the aggregate matching function} \\
\bigskip

Let's recall that the job finding rate $F_{t}=\frac{H_{t}}{U{t}}$\\
\bigskip

The state space representation in this case is the following:

\begin{align}
F_{t} &= M_{t} + \rho (V_{t}-U_{t})+\epsilon_{t} \text{     } \epsilon_{t} \sim N(0,R) \\
M_{t}&=M_{t-1}+\nu_{t} \text{     } \nu_{t} \sim N(0,Q)
\end{align}
where also, $Cov(\epsilon_{t},\nu_{t})=0$ and $(V_{t}-U_{t})$ the labor market tightness.

\end{frame}

\begin{frame}

\frametitle{First Step:}
\textbf{What to do} \\
\bigskip

1 - In the $Main\_File.m$ set the initial values for the three estimated parameters (Q, R, $\rho$). (Think of OLS as starting point).\\ \bigskip

2 -  Next, in $Loglikelihood.m$ specify the estimated parameters in terms of the notation used in previous section. That is, determine the values of Q, R, C, F, M and A given the input of the minimization routine. \\ \bigskip

3 - Finally, complete the Kalman filter recursions in $KalmanFilter.m$ Then run the program by commenting out the rest of the code and see what happened to match efficiency during the crisis.
\end{frame}


\miniframesoff
\section{References}

\miniframeson

\begin{frame}

\frametitle{References}

Den Haan, W. (2019). ''Bayesian Estimation'', Tools for Macroeconomists

Petrongolo, B. , C. Pissarides (2001). ''Looking into the black box: A survey of the
matching function", Journal of Economic Literature, 39, 390-431. \\ \bigskip

Sedlacek, P. (2019). ''Kalman filter and Maximum Likelihood'', Tools for Macroeconomists\\ \bigskip

\end{frame}
\end{document}